% 第2回　問いの立て方（記入例・模範解答）
\session{2}{1}{10月7日（水）}{問いの立て方}{AIとブレスト}{}

\goals{
  \item テーマとリサーチ・クエスチョンの違いを説明できる
  \item 良い問いの条件（調べて答えられる／範囲が適切／答える意義がある）を挙げられる
  \item AIとのブレストで視点を広げつつ、最後は自分で選ぶ
}

\work{講義メモ}{テーマと問いの違い}
\begin{wstab}{colspec={Q[l,wd=40mm,m] X[l,m]},row{1-Z}={ht=9mm}}
  \hc{例：テーマ} & 地方の自動車販売 \\
  \hc{例：問い} & なぜ若い世代の来店は減っているのか \\
  \hc{自分のテーマ} & \af{若い世代と自動車ディーラー} \\
  \hc{問いにすると} & \af{名古屋都市圏の20代は、なぜディーラーに足を運ばなくなっているのか} \\
  \hc{良い問いの3条件} & ①\,\af{調べて答えられる}\hspace{8mm} ②\,\af{範囲が広すぎず狭すぎない}\hspace{8mm} ③\,\af{答える意義がある} \\
\end{wstab}

\work[60mm]{ワーク1}{AIで視点を広げる}
\begin{promptbox}[視点を広げる]
私は経済学部の2年生です。「（あなたのテーマ）」に関心があります。このテーマについて、経済学、経営学、社会学、心理学、都市計画、テクノロジーの6つの視点から、それぞれ2つずつ研究の問いを提案してください。各問いは「なぜ〜なのか」「〜は〜にどのように影響するのか」の形で、1行ずつ箇条書きにしてください。
\end{promptbox}
\instr{AIの候補のうち、気になったものを視点ごとに書き写す（言い換えてよい）。}
\begin{wstabh}{colspec={Q[l,wd=22mm,m] X[l,m]},row{2-Z}={ht=10mm}}
  視点 & 気になった問いの候補 \\
  経済学 & \af{若者の所得と維持費の変化は、車の保有にどう影響しているのか} \\
  経営学 & \af{ディーラーの販売手法は、若い顧客の来店にどう対応しているのか} \\
  社会学 & \af{「車を持つこと」の意味は、世代によってどう変わったのか} \\
  心理学 & \af{なぜ若い客は、販売員と対面で話すことを負担に感じるのか} \\
  都市計画 & \af{公共交通の便利さは、都市部の若者の車離れにどう影響しているのか} \\
  テクノロジー & \af{オンラインでの情報収集は、ディーラーへの来店のしかたをどう変えたのか} \\
\end{wstabh}

\work[60mm]{ワーク2}{3条件で採点し、上位3つに絞る}
\instr{各条件を 3（十分）・2（やや）・1（不十分）で採点する。必要ならAIに評価と狭める案を求める。}
\begin{promptbox}[問いを絞る]
次の問いについて、「大学2年生が半年で調べて答えられるか」「範囲は適切か」「答えることに意義があるか」の3点から評価し、必要なら狭める案を示してください。問い：（候補の問い）
\end{promptbox}
\begin{wstabh}{colspec={X[l,m] Q[c,wd=17mm,m] Q[c,wd=17mm,m] Q[c,wd=17mm,m] Q[c,wd=12mm,m] Q[c,wd=10mm,m]},row{2-Z}={ht=10mm}}
  問いの候補 & 調べて\newline 答えられる & 範囲が\newline 適切 & 答える\newline 意義 & 合計 & 順位 \\
  \af{オンラインの情報収集は来店のしかたをどう変えたか} & \ans{3} & \ans{2} & \ans{3} & \ans{8} & \ans{1} \\
  \af{なぜ若い客は販売員との対話を負担に感じるのか} & \ans{2} & \ans{3} & \ans{3} & \ans{8} & \ans{2} \\
  \af{公共交通の便利さと若者の車離れ} & \ans{3} & \ans{2} & \ans{2} & \ans{7} & \ans{3} \\
  \af{「車を持つこと」の意味の世代差} & \ans{1} & \ans{1} & \ans{3} & \ans{5} & \\
  \af{若者の所得と維持費の変化} & \ans{2} & \ans{1} & \ans{2} & \ans{5} & \\
\end{wstabh}

\work{ペアワーク}{問いを説明し合う}
\atwoboxes{相手からの質問（「それはどうやって調べるの」など）}{自分の答え・気づいたこと}{30mm}
{「来店のしかたが変わった」と、どうやって確かめるの？ 来店の回数？ 来店の目的？}
{来店の回数と、来店前に何を調べたかを、同世代へのアンケートと公的調査で見たい。「変わった」の中身を決めないと調べられないと気づいた。}

\work{決定}{仮のリサーチ・クエスチョン}
\abox{18mm}{名古屋都市圏の20代にとって、オンラインでの情報収集は、自動車ディーラーへの来店のしかたをどう変えたのか。}

\begin{nexttime}
  \item 仮のリサーチ・クエスチョンを1〜2文で書き、選んだ理由を200字程度でまとめる
  \item 使ったプロンプトと、AIの案のうち採用しなかったものとその理由をメモする（下のAI活用記録）
\end{nexttime}
\abox[選んだ理由（200字程度）]{40mm}{実家の近くのディーラーがいつも空いている一方で、同世代は車の情報を動画やサイトで詳しく調べている。若い世代は車に関心がないのではなく、関心の向け方や店への行き方が変わったのではないかと考えた。この問いは、公的調査やアンケートで確かめられる範囲にあり、来店客の減少に悩む販売店にとっても答える意義がある。後半の店舗設計の課題ともつながるので、半年を通して深められると考えた。}
\aimemoA{問いの候補出し（6つの視点）と、候補の採点}
{「6つの視点から問いを2つずつ」「3条件で評価し、狭める案を示して」}
{12の候補から6つを採用。所得と維持費の問いは、統計の扱いが難しく範囲が広いため不採用。採点はAIの評価を参考にしたが、意義の点数は自分で付け直した。}
{AIが挙げた「若者の車離れ」が実際に指摘されているかを、白書や業界団体の調査があるかで確かめた（次回に詳しく調べる）}

\begin{point}
  \item テーマ（話題の範囲）と問い（答えを求める疑問文）の違いを、自分のテーマで書き分けられているかを見る。疑問文になっていない「〜について」は問いではない。
  \item 良い問いの3条件は、授業で示した「調べて答えられる」「範囲が適切」「答える意義がある」の3つが模範解答。
  \item 採点表は数値の妥当性より、点数の理由を口頭で説明できるかが大切。AIの評価をそのまま写していないかを、ペアワークでの説明から確かめる。
  \item AI活用記録に「採用しなかった案と理由」が書けているかを必ず見る。第5回のレポートで求める記録の練習になっている。
\end{point}
